\begin{proposition}\hypertarget{tag-01Q7}{}\label{fam:src:B.1.1}\label{fam:st:B.1.1}
\StatementRecord{fam:src:B.1.1}{proved}{01Q7}{\SeriesRef{fam:src:A.3.1}}

Let $E_\tau=\mathbb C/(\mathbb Z+\tau\mathbb Z)$, $P=[\mu]$, and $L_\mu=\mathcal O_{E_\tau}(P-O)$. The holomorphic map
\[
 [\xi,\eta]\longmapsto[\xi,e^{2\pi i\eta}]
\]
identifies $\mathbb C^2/\Pi_c\Lambda$ with the quotient of $L_\mu^\times$ by the lift
$A(\xi,W)=(\xi+6\mu,e^{2\pi i\beta_c}W)$. Engel's generator, in Definition 2.3 of the attached paper, is its inverse, with translation $-6P$ and multiplier $e^{-2\pi i\beta_c}$. The section $P=[\mu]$ is a generator of the Mordell--Weil group of the stated pointed Weierstrass surface.
\end{proposition}


\begin{proof}
\begin{enumerate}[label=\arabic*.,leftmargin=*]
\item First quotient by columns two, three and four. Exponentiation kills column four; columns three and two then give
\[
 (\xi,W)\sim(\xi+1,W)\sim(\xi+\tau,e^{2\pi i\mu}W).
\]
For the odd theta function with periods $1,\tau$, the quotient $\theta(\xi-\mu)/\theta(\xi)$ is a meromorphic section of this line bundle: its factors are $1,e^{2\pi i\mu}$ and its divisor is $P-O$. Sending its divisor section to this quotient is the asserted line-bundle isomorphism. Quotienting the remaining first column gives exactly $A$. On sufficiently small simply connected charts the logarithm of $W$ is a holomorphic inverse; changes of logarithm are column-four periods. Thus the displayed map is a biholomorphism of the two quotient complex spaces, with its inverse explicitly defined by local logarithms. Choosing $A^{-1}$ changes the generator of the same infinite cyclic group and gives precisely Engel's declared inverse-translation convention.

\item For clarity the elliptic surface in this dictionary is the displayed cubic, rather than only a curve with the same $j$-value. Let $f_{\rm sc}$ be the branch satisfying
\[
 f_{\rm sc}^4=\frac{12t^3(t-1)}{g_2(\tau)},\qquad
 f_{\rm sc}^6=-\frac{8t^4(t-1)^2}{g_3(\tau)},\qquad
 f_{\rm sc}/t\longrightarrow i\sqrt3/\pi.
\]
The identities $g_2^3/(g_2^3-27g_3^2)=t$ make these equations compatible. Their zero and pole orders give a holomorphic branch on the period cover; the indicated cusp sign fixes it. The map $\xi\mapsto(f_{\rm sc}^2\wp_\tau(\xi),\frac12f_{\rm sc}^3\wp'_\tau(\xi))$ is a pointed isomorphism to
$y^2=x^3-3t^3(t-1)x+2t^4(t-1)^2$. Its factors are $f_{\rm sc}(g_1z)=-f_{\rm sc}/\tau$, $f_{\rm sc}(g_2z)=f_{\rm sc}/\tau$; therefore this isomorphism descends over $B^\circ$.

\item We first verify that $P=[\mu]$ gives a rational section of the resolved elliptic surface. On $B^\circ$ its cubic coordinates descend by the displayed scale and affine laws. At the two finite centres, $\mu_j$ is not a lattice point, while $f_{\rm sc}$ has orders two and one in the root parameter. Its cubic coordinates are therefore holomorphic and invariant, hence holomorphic in the base coordinate. At the cusp $\mu$ is holomorphic and bounded, and the convergent multiplicative expansion of $\wp$ gives meromorphic cubic coordinates with finite pole orders. Even when its exponential meets one, the pole is controlled by its finite vanishing order; it cannot be identically one because the two finite values of $\mu$ differ. The coordinates are thus meromorphic on $\mathbb P^1$ and rational. The closure of its graph on the proper resolved surface is an irreducible curve; its normalization maps finitely with degree one to the smooth base, so is the base itself. This gives the unique section of that surface extending $P$.

\item We also verify that the specified section is a Mordell--Weil generator using its full logarithm. A rational section of this cubic pulls back to a holomorphic section of the good elliptic family at both finite covering points: multiply its projective coordinates by a common power to remove poles, and use the nonsingular limiting Weierstrass equation. At the cusp put $v=1/t$, $X=v^2x,Y=v^3y$. Its equation is $Y^2=X^3-3(1-v)X+2(1-v)^2$. A section cannot pass through the node $(1,0)$: substitution of $X=1+O(v)$, $Y=O(v)$ leaves the coefficient $+1$ of $v$ on the left-hand side of the defining equation. At every other central point the multiplicative uniformization is a local isomorphism. Explicitly, at $v=0$ it is
\[
 X=1+\frac{12U}{(U-1)^2},\qquad
 Y=\frac{12\sqrt3 U(U+1)}{(U-1)^3},\qquad U=e^{2\pi i\xi},
\]
with inverse $U=(Y/(X-1)+\sqrt3)/(Y/(X-1)-\sqrt3)$ away from $O$; a local parameter at $O$ is $X/Y$. The convergent lattice expansion of $\wp$ gives the same holomorphic map for small $v$. Hence a section has a holomorphic unit $U(v)$ and a bounded logarithm at the cusp. Lifting the section on the simply connected cover gives a global logarithm $\nu$, with this cusp branch. Its only remaining ambiguity is an integer, since a nonzero multiple of $\tau$ is unbounded.

\item Its two affine laws have integral defects:
\[
 \nu(g_1z)=-\nu/\tau+\alpha_1+\beta_1\tau(g_1z),\quad
 \nu(g_2z)=\nu/\tau+\alpha_2+\beta_2\tau(g_2z).
\]
Composing them and using cusp invariance gives $\alpha_2=-\beta_1$, $\beta_2=\alpha_1+\beta_1$. Put $n=-\alpha_1-2\beta_1$ and $a=\alpha_1+\beta_1$. Substitution shows that $n\mu+a$ has exactly these defects. Their difference is a global section of the first-coordinate sheaf $\mathcal O_B(-p_0)$ from \ref{fam:src:A.3.1}, so is zero. Every rational section is therefore $nP$. The same cusp argument and the distinct fixed values of $\mu$ show that $P$ has infinite order. This proves $\operatorname{MW}(S)=\mathbb ZP$.
\end{enumerate}
\end{proof}